\documentclass[10pt]{article}
\usepackage[margin=12mm]{geometry}
\usepackage{graphicx}
\usepackage{caption}
\pagestyle{empty}
\captionsetup{font=small,labelfont=bf}

\begin{document}
\begin{figure}[p]
\centering
\includegraphics[width=0.95\textwidth]{figS3A_ce_tracking_replicate_fronts.pdf}\\[2mm]
\includegraphics[width=0.95\textwidth]{figS3B_ce_tracking_replicate_metrics.pdf}

\caption{\textbf{The terminal-cell Pareto result is globally robust across
three \textit{C. elegans} tracking replicates but retains local subtree
dependence.}
%
The predefined stage-matched cutoffs differ because embryonic development
proceeds at different rates across tracks (embryo 1, $T=255$; embryo 2,
$T=247$; embryo 3, $T=226$). The respective tracks contain 299, 301, and 309
usable terminal parent--child edges. All comparisons use the strict
intersection of 279 edges (93.3\%, 92.7\%, and 90.3\% of the available edges,
respectively), so terminal identities, natural parentage, and the top-20
protein representation are identical across replicates.
%
\textbf{(A)} Pareto fronts for the matched all-terminal set. Travel and
cell-state distances are separately expressed in standard-deviation units
from each replicate's first-cousin-shuffle null and translated so its natural
lineage is at $(0,0)$. Color and line style identify tracking replicate. The
black cross marks the shared natural assignment, blue diamonds mark each
front's maximum-edge-retention assignment, gold points show 1,000
first-cousin shuffles, and gold plus signs mark the corresponding exact null
means. The three fronts remain closely aligned under the matched analysis.
%
\textbf{(B)} Replicate-level canonical position $u_L$, endpoint-normalized
natural-lineage-to-front distance $d_{LP}$, and maximum edge retention for P0,
AB, ABa, ABp, and P1. As in Figure 5, $u_L$ and $d_{LP}$ use the closest
sampled Pareto assignment in endpoint-normalized geometry; interpolated
segments are not treated as attainable lineages. At each subtree position,
three semi-transparent, vertically aligned horizontal marks show the replicate
values; mark color matches the replicate encoding in panel A, and coincident
values blend visibly. The P0 values vary modestly, whereas ABa is more
replicate-sensitive: the embryo-3
natural assignment lies on the sampled front and reaches unit maximum edge
retention. The objectives remain separate and are not combined into an
efficiency score.}
\label{fig:ce_tracking_robustness}
\end{figure}
\end{document}
